\section{Cartan genericity and a residual linear system}\label{sec:linear}
Throughout this section $r\geq3$, and the eigenvalues $\mu_i$ are those
fixed in Lemma~\ref{lem:units}. All parameter spaces are over $\C$.

\begin{lemma}\label{lem:P}
There is a nonempty Zariski open set of $Q\in\SL_r$ on which both of the
following hold:
\begin{enumerate}[label=(P\arabic*),leftmargin=3em]
\item For every $i\ne j$, the map $\hc\to\C^2$,
$K\mapsto(K_{ij},K_{ji})$, has rank two.
\item The torus $\TC$ contains a regular matrix $N$ whose off-diagonal
entries are all nonzero and such that, for every three distinct indices,
\begin{equation}\label{eq:triangle-generic}
 N_{ij}N_{jk}N_{ki}\ne N_{ik}N_{kj}N_{ji}.
\end{equation}
\end{enumerate}
\end{lemma}
\begin{proof}
For (P1) at the pair $(1,2)$, take a matrix $H$ whose only nonzero
off-diagonal entries are
\[
 H_{12}=H_{21}=H_{23}=H_{31}=1;
\]
\emph{all other off-diagonal entries are zero}. Give it generic diagonal
entries. The projections of $H$ and $H^2$ to entries $(1,2),(2,1)$ are
\[
 (1,1),\qquad(H_{11}+H_{22},\ H_{11}+H_{22}+1),
\]
so their determinant is one. Generic diagonal entries make $H$ regular
semisimple (one may first take distinct widely separated real diagonal
entries). Subtract the scalar part of $H$; this does not change the
independence, since the change in $H^2$ is a scalar plus a multiple of $H$.
The two traceless polynomial expressions belong to the eigen-Cartan of
$H$. Choosing an eigenbasis and rescaling one column makes its matrix
$Q$ have determinant one. The rank-two condition is an open condition on
$Q$, since it is the nonvanishing of at least one $2\times2$ minor of a
fixed basis of $Q\h Q^{-1}$. Relabeling the witness treats every pair.
The finitely many nonempty opens intersect in the irreducible variety
$\SL_r$.

For (P2), the product of the off-diagonal coordinates, the finitely many
triangle differences in \eqref{eq:triangle-generic}, the determinant,
and the discriminant is a nonzero polynomial on the full matrix space:
each factor is a nonzero polynomial, and that polynomial ring is a domain.
Choose a matrix where it is nonzero and scale it to determinant one.
Its eigen-Cartan is a witness. More explicitly, substitute
$N=Q\diag(t_1,\ldots,t_r)Q^{-1}$, $\prod t_i=1$, in the product of the
required factors. This is a Laurent polynomial in $t_1,\ldots,t_{r-1}$
whose coefficients are regular in $Q$. Its not being identically zero is
a nonempty open condition on $Q$. It is equivalent to the existence of
a suitable $N$. Intersect with (P1).
\end{proof}

\begin{lemma}[Principal-minor differential rigidity]\label{lem:minors}
Suppose a matrix $M\in M_r(\C)$ has no zero off-diagonal entries and
satisfies \eqref{eq:triangle-generic} with $M$ in place of $N$. If a
variation $E\in M_r(\C)$ makes every principal minor stationary, then
$E=[h_0,M]$ for a diagonal matrix $h_0$.
\end{lemma}
\begin{proof}
The one-by-one minors give $E_{ii}=0$. Set $x_{ij}=E_{ij}/M_{ij}$ for
$i\ne j$. The two-by-two minors give $x_{ij}+x_{ji}=0$.
Differentiate the three-by-three determinant by the permutation formula.
Its diagonal term vanishes because $E_{ii}=0$; its transposition terms
vanish by the pair identities. The remaining two oriented cycles give
\[
 (M_{ij}M_{jk}M_{ki}-M_{ik}M_{kj}M_{ji})
                 (x_{ij}+x_{jk}+x_{ki})=0.
\]
The first factor is nonzero, so all oriented triangle sums vanish.
Set $h_1=0$ and $h_i=x_{i1}$ for $i\ne1$. Antisymmetry and the triangle
with reference index $1$ give $x_{ij}=h_i-h_j$ for every $i\ne j$,
including the cases involving index $1$. Thus
$E_{ij}=(h_i-h_j)M_{ij}$, including the diagonal entries.
\end{proof}

This is an infinitesimal form of classical principal-minor rigidity under
diagonal similarity. Loewy's global theorem \cite{Loewy} and the
pair-product and triangle-cycle identities of Mantelos
\cite[Sections 3--4 and Lemma 6.2]{Mantelos} are antecedents.
We give the elementary proof because its exact differential hypotheses
are what is needed below; no new contribution is claimed for this lemma
alone. In rank three the three-by-three minor is the full determinant;
in both applications its derivative vanishes automatically.

\begin{proposition}\label{prop:residual}
There is a nonempty Zariski open subset of $T\times\SL_r$ such that,
for $C=QC_\mu Q^{-1}$, the only $S\in\slr$ satisfying
\begin{equation}\label{eq:residual}
 (\Ad_A-\Ad_C)S=0,\qquad (I-\Ad_C)S\in\h+\hc
\end{equation}
is zero, and $\h\cap\hc=0$.
\end{proposition}
\begin{proof}
Choose a real orthogonal matrix $Q$ of determinant one with every entry
nonzero. Such matrices exist: each individual coordinate is nonzero on
an open dense subset of the real analytic manifold $\mathrm{SO}(r)$,
and a finite intersection is nonempty. Every member of $\hc$ is then a
complex symmetric matrix. Moreover $\h\cap\hc=0$: if
$\diag(k)Q=Q\diag(\lambda)$, nonzero $Q_{ij}$ implies
$k_i=\lambda_j$ for all $i,j$, so both diagonal matrices are scalar, and
tracelessness makes them zero.

Choose a real traceless diagonal $H$ with distinct entries, and put
$A(t)=\exp(tH)$ for real $t$. Suppose there were nonzero solutions of
\eqref{eq:residual} for a sequence $t\to0$, $t\ne0$. Normalize the
Hermitian norms of $S_t$ to one and pass to a convergent subsequence
$S_t\to K$. The first equation gives $K\in\hc$. It also gives
\[
 \frac{(I-\Ad_C)S_t}{t}
 =\frac{(I-\Ad_{A(t)})S_t}{t}\longrightarrow-[H,K].
\]
The left-hand side lies in the fixed closed subspace $\h+\hc$, hence
$[H,K]$ is complex symmetric. Since $H$ and $K$ are symmetric, their
commutator is skew-symmetric. Therefore $[H,K]=0$. Distinct entries of
$H$ make $K$ diagonal, whence $K\in\h\cap\hc=0$, contradicting its
norm one. Thus for all sufficiently small nonzero $t$ the system is
injective.

For global openness use the augmented linear map
\begin{equation}\label{eq:augmented}
\begin{split}
 \slr\oplus\h\oplus\h&\longrightarrow\slr\oplus\slr,\\
 (S,\alpha,\beta)&\longmapsto
 \bigl((\Ad_A-\Ad_C)S,\ (I-\Ad_C)S-\alpha-Q\beta Q^{-1}\bigr).
\end{split}
\end{equation}
It is injective exactly when \eqref{eq:residual} has no nonzero solution
and $\h\cap\hc=0$. Indeed a vector in its kernel gives the residual
system; after $S=0$, the remaining equality is exactly an intersection
of the two Cartans. Conversely each residual solution admits such
$\alpha,\beta$. All entries of \eqref{eq:augmented} are regular functions
on $T\times\SL_r$. Injectivity is a maximal-minor condition, and the
witness just constructed proves its nonemptiness.
\end{proof}

\begin{lemma}\label{lem:all-open}
There is a nonempty open set of triples $(A,B,Q)\in T\times T\times\SL_r$
satisfying (P1), (P2), Proposition~\ref{prop:residual}, and the following
additional conditions: $B$ has distinct diagonal entries; every entry of
$Q$ and $C=QC_\mu Q^{-1}$ is nonzero; and for each fixed cardinality
$1\leq j<r$, the products
\[
 a_I=\prod_{i\in I}a_i\qquad(|I|=j),\quad A=\diag(a_1,\ldots,a_r),
\]
are distinct for distinct subsets $I$.
\end{lemma}
\begin{proof}
Each condition is nonempty and open, and the parameter space is
irreducible. For the subset products, $\mathbf1_I-\mathbf1_J$ is not a
multiple of $(1,\ldots,1)$ when $I\ne J$ have equal size, so the forbidden
character equality is proper. Nonzero entries of $Q$ and regularity of
$B$ are elementary open conditions. Each off-diagonal entry of a matrix
in the regular conjugacy class of $C_\mu$ is nonzero somewhere: conjugate
a diagonal matrix by an elementary matrix and use two distinct eigenvalues.
Each diagonal entry is nonzero somewhere, since the eigenvalues themselves
are nonzero. Their individual nonempty opens therefore intersect.
The other conditions are covered by the preceding results.
\end{proof}
